\section{Adiabatik jarayon}


\begin{enumerate}

\item Qaysi qonun tabiatdagi jarayonlarning yo‘nalishini ko‘rsatadi?
\begin{enumerate}[label=\Alph*)]
\item massaning saqlanish qonuni
\item termodinamikaning I qonuni
\item termodinamikaning II qonuni
\item energiyaning saqlanish qonuni
\end{enumerate}

\item Issiqlik jihatdan izolyatsiyalangan sistemadagi gaz siqilsa, uning temperaturasi qanday o‘zgaradi?
\begin{enumerate}[label=\Alph*)]
\item ortadi
\item pasayadi
\item o‘zgarmaydi
\item avval ortadi so‘ng pasayadi
\end{enumerate}

\item Adiabatik kengayishda gazning ichki energiyasi qanday o‘zgaradi?
\begin{enumerate}[label=\Alph*)]
\item $\Delta U < 0$
\item $\Delta U = 0$
\item $\Delta U > 0$
\item o‘zgarmaydi
\end{enumerate}

\item Termodinamikaning 1–qonuni adiabatik jarayon uchun qanday ko‘rinishda yoziladi?
\begin{enumerate}[label=\Alph*)]
\item $Q = \Delta U + A$
\item $Q = \Delta U$
\item $A = Q$
\item $\Delta U + A = 0$
\end{enumerate}

\item Ideal gazni adiabatik tarzda siqib, ish bajarildi. Bunda gaz molekulalarining kinetik va potensial energiyalari qanday o‘zgaradi?
\begin{enumerate}[label=\Alph*)]
\item kinetik energiya ortadi, potensial energiya kamayadi.
\item kinetik energiya kamayadi, potensial energiya o‘zgarmaydi.
\item kinetik va potensial energiyalar ortadi.
\item kinetik energiya ortadi, potensial energiya o‘zgarmaydi.
\end{enumerate}

\item Issiqlik jihatdan izolyatsiyalangan sistemadagi gaz siqilsa, uning temperaturasi qanday o‘zgaradi?
\begin{enumerate}[label=\Alph*)]
\item ortishi ham, pasayishi ham mumkin
\item pasayadi
\item ortadi
\item o‘zgarmaydi
\end{enumerate}

\item Gaz adiabat ravishda siqilmoqda. Bunda gazning $T$ temperaturasi va $P$ bosimi qanday o‘zgaradi?
\begin{enumerate}[label=\Alph*)]
\item $T$ ortadi, $P$ kamayadi
\item $T$ kamayadi, $P$ ortadi
\item $T$ va $P$ ortadi
\item $T$ va $P$ kamayadi
\end{enumerate}

\item Adiabatik jarayonda gazning ichki energiyasi 5 kJ kamaydi. Gaz necha kJ ish bajardi?
\begin{enumerate}[label=\Alph*)]
\item 5
\item 2,5
\item 10
\item 0,5
\end{enumerate}

\item Porshen ostida 10 mol kislorod bor. U adiabatik ravishda kengayib $10^\circ C$ ga sovidi. U tashqi kuchlarga qarshi qancha ish bajaradi (J)?
\begin{enumerate}[label=\Alph*)]
\item 2820,5
\item 2659,5
\item 1246,5
\item 3112,5
\end{enumerate}

\item Gaz ustida tashqi kuchlar manfiy ish bajarsa, uning zichligi qanday o‘zgaradi?
\begin{enumerate}[label=\Alph*)]
\item kamayadi
\item ortadi
\item o‘zgarmaydi
\item kamayishi ham ortishi ham mumkin
\end{enumerate}

\item Gaz ustida tashqi kuchlar musbat ish bajarganda uning zichligi qanday o‘zgaradi?
\begin{enumerate}[label=\Alph*)]
\item ortishi ham kamayishi ham mumkin
\item ortadi
\item o‘zgarmaydi
\item kamayadi
\end{enumerate}

\item Adiabatik jarayonda gaz musbat ish bajarsa, uning zichligi qanday o‘zgaradi?
\begin{enumerate}[label=\Alph*)]
\item o‘zgarmaydi
\item kamayadi
\item ortadi
\item ortishi ham kamayishi ham mumkin
\end{enumerate}

\item Ideal gaz siqildi, bunda adiabatik jarayon yuz berdi. Gazning temperaturasi qanday o‘zgaradi?
\begin{enumerate}[label=\Alph*)]
\item kamayadi
\item ortadi
\item ortishi ham kamayishi ham mumkin
\item o‘zgarmaydi
\end{enumerate}

\item Ideal gaz kengaydi, bunda adiabatik jarayon yuz berdi. Gazning temperaturasi qanday o‘zgaradi?
\begin{enumerate}[label=\Alph*)]
\item kamayadi
\item ortadi
\item ortishi ham kamayishi ham mumkin
\item o‘zgarmaydi
\end{enumerate}

\item Ideal gazning ma'lum bir miqdori $V_1$ hajmdan $V_2$ hajmgacha uch marta siqildi. Keltirilgan grafikdagi qaysi egri chiziq eng sekin siqilishga mos keladi?
\begin{enumerate}[label=\Alph*)]
\item 1
\item 2
\item 3
\item TJY
\end{enumerate}

\item Ideal gazning ma'lum bir miqdori $V_1$ hajmdan $V_2$ hajmgacha uch marta siqildi. Keltirilgan grafikdagi qaysi egri chiziq 1 sekund davomidagi siqilishga mos keladi?
\begin{enumerate}[label=\Alph*)]
\item 2
\item 1
\item 3
\item TJY
\end{enumerate}

\end{enumerate}


\begin{table}[h!]
\centering
\begin{tabular}{|c|c|c|c|c|c|c|c|c|c|}
\hline
№ & J. & № & J. & № & J. & № & J. & № & J. \\ \hline
1 & C & 5 & D & 9 & C & 13 & B & & \\ \hline
2 & A & 6 & C & 10 & A & 14 & A & & \\ \hline
3 & A & 7 & C & 11 & B & 15 & A & & \\ \hline
4 & D & 8 & A & 12 & B & 16 & A & & \\ \hline
\end{tabular}
\end{table}